\section*{Hoofdstuk \ref{ch:b1:work-and-energy} --- Arbeid, energie en potentiële energie}

\begin{solution}{exo:b1:work-and-energy:1}
Hefkracht $+mgh = 20 \times 9.81 \times 3.0 = \qty{589}{J}$; gewicht
$-\qty{589}{J}$; vermogen $589/5.0 = \qty{118}{W}$. Wrijving: $-\mu_dmgd =
-0.30 \times 196 \times 3.0 = \qty{-177}{J}$.
\end{solution}

\begin{solution}{exo:b1:work-and-energy:2}
$E_k = \tfrac12 \times 1000 \times 27.8^2 = \qty{386}{kJ}$; $d = E_k/F =
\qty{55}{m}$; $v$ verdubbelen viervoudigt $E_k$ en $d$: \qty{220}{m}.
\end{solution}

\begin{solution}{exo:b1:work-and-energy:3}
$E_p = \tfrac12 \times 500 \times 0.10^2 = \qty{2.5}{J}$; $v = \sqrt{2 \times
2.5/0.050} = \qty{10}{m/s}$.
\end{solution}

\begin{solution}{exo:b1:work-and-energy:4}
$\tfrac12 mv^2 = mg\ell(1 - \cos 60^\circ)$: $v = \sqrt{g\ell} = \qty{3.1}{m/s}$.
Spankracht: $T - mg = mv^2/\ell = mg$, dus $T = 2mg$.
\end{solution}

\begin{solution}{exo:b1:work-and-energy:5}
De wrijving verricht $-\mu_dmg\cos\alpha\,(h/\sin\alpha)$, dus
\[
\tfrac12 mv^2 = mgh(1 - \mu_d\cot\alpha) = mgh(1 - 0.35),
\]
$v = \sqrt{2 \times 9.81 \times 2.0 \times 0.65} = \qty{5.1}{m/s}$; $35\%$ verloren.
\end{solution}

\begin{solution}{exo:b1:work-and-energy:6}
$E_m = \tfrac12 mv^2 - GMm/R \geq 0$ opdat het lichaam het oneindige bereikt
(waar $E_p = 0$ en $v \to 0$): $v_e = \sqrt{2GM/R}$, dus
\[
v_e = \sqrt{\frac{2 \times 6.67\times10^{-11} \times 5.97\times10^{24}}{6.37\times10^6}}
= \qty{11.2}{km/s}.
\]
Bij dezelfde dichtheid is $M \propto R^3$ en dus $v_e \propto R$: \qty{22.4}{km/s}.
\end{solution}

\begin{solution}{exo:b1:work-and-energy:7}
Laagste punt: $mg(\ell_0 + x) = \tfrac12 kx^2$, $25x^2 - 687x - 13734 = 0$,
$x = \qty{41}{m}$; totale valhoogte \qty{61}{m}; $T_{\max} = kx = \qty{2.0}{kN}
= 3.0\,mg$ (netto \qty{2}{g} omhoog). De snelheid is het grootst waar de netto kracht
verdwijnt, $kx = mg$: $x = \qty{14}{m}$, \qty{34}{m} onder het sprongpunt.
\end{solution}

\begin{solution}{exo:b1:work-and-energy:8}
Zwaartekracht: $mgv\sin\alpha = 80 \times 9.81 \times 5.0 \times 0.05 = \qty{196}{W}$;
weerstand: $\tfrac12\rho C_xSv^3 = 0.5 \times 1.2 \times 0.50 \times 125 =
\qty{38}{W}$; totaal \qty{234}{W}.
\end{solution}

\begin{solution}{exo:b1:work-and-energy:9}
$E_p' = 4E_0(x^3/a^4 - x/a^2) = 0$: $x = 0, \pm a$. $E_p'' = E_0(12x^2/a^4 -
4/a^2)$: $-4E_0/a^2$ in $0$ (labiel maximum, $E_p = 0$), $+8E_0/a^2$
in $\pm a$ (stabiele minima, $E_p = -E_0$): putten met diepte $E_0$;
$\omega_0^2 = 8E_0/ma^2$.
\end{solution}

\begin{solution}{exo:b1:work-and-energy:10}
$\tfrac12 m\dot x^2 + \tfrac12 kx^2 = E$: een ellips met halve assen $x_{\max}
= \sqrt{2E/k}$ en $\dot x_{\max} = \sqrt{2E/m}$. De beweging is harmonisch
met periode $2\pi\sqrt{m/k}$, wat de \omterm{def:b1:wave-propagation:signal}{amplitude} ook is (isochronie), zodat
elke ellips dezelfde tijd kost.
\end{solution}

\begin{solution}{exo:b1:work-and-energy:11}
Energie: $v^2 = 2gR(1 - \cos\theta)$. Radiaal: $mg\cos\theta - N = mv^2/R$,
dus $N = mg(3\cos\theta - 2)$, nul bij $\cos\theta = 2/3$ ($\theta =
48^\circ$): hoogte $2R/3$, snelheid $\sqrt{2gR/3}$.
\end{solution}

\begin{solution}{exo:b1:work-and-energy:12}
$\tfrac12 mv^2 = \mu_dmgd$: $\mu_d = 625/(2 \times 9.81 \times 80) = 0.40$.
Warmte in het contact band--wegdek (en een remspoor). Rollende wielen gebruiken
statische wrijving, $\mu_s > \mu_d$: een grotere remkracht, een kortere
afstand, en de besturing blijft behouden.
\end{solution}

\begin{solution}{pb:b1:work-and-energy:1}
\textbf{1.} $r \to 0$: $+\infty$ (de term $r^{-12}$); $r \to \infty$:
$0^-$. Een steile wand, een put, een lange staart.

\textbf{2.} $E_p' = \epsilon[-12\sigma^{12}/r^{13} + 12\sigma^6/r^7] = 0$ precies als
$(\sigma/r)^6 = 1$: $r = \sigma$, $E_p = \epsilon(1 - 2) = -\epsilon$.

\textbf{3.} $(\sigma/r)^{12} = u^2$: $E_p = \epsilon(u^2 - 2u) = \epsilon[(u - 1)^2
- 1]$, minimaal bij $u = 1$ ($r = \sigma$) met waarde $-\epsilon$.

\textbf{4.} $F = -E_p' = (12\epsilon/\sigma)[(\sigma/r)^{13} - (\sigma/r)^7]$:
positief (afstotend) voor $r < \sigma$, negatief (aantrekkend) daarbuiten.

\textbf{5.} $(1/1.2)^6 = 0.335$, $(1/1.2)^{12} = 0.112$: $E_p = \epsilon(0.112
- 0.670) = -0.56\epsilon = \qty{-2.5}{eV}$. $(1/1.2)^{13} = 0.0935$,
$(1/1.2)^7 = 0.279$: $F = (12\epsilon/\sigma)(-0.186) = -2.2\epsilon/\sigma =
-2.2 \times 7.0\times10^{-19}/1.27\times10^{-10} = \qty{-12}{nN}$.

\textbf{6.} $\epsilon = \qty{4.4}{eV} = 4.4 \times 96.5 = \qty{425}{kJ/mol}$,
dicht bij de gemeten \qty{431}{kJ/mol}.

\textbf{7.} $E_p'' = \epsilon[156\sigma^{12}/r^{14} - 84\sigma^6/r^8]$; bij
$\sigma$: $72\epsilon/\sigma^2 > 0$, dus stabiel.

\textbf{8.} $E_p \approx -\epsilon + \tfrac12(72\epsilon/\sigma^2)(r - \sigma)^2$:
$k = 72\epsilon/\sigma^2$.

\textbf{9.} $k = 72 \times 7.0\times10^{-19}/(1.27\times10^{-10})^2 =
\qty{3.1e3}{N/m}$.

\textbf{10.} $\omega_0 = \sqrt{k/m} = \sqrt{3.1\times10^3/1.67\times10^{-27}} =
\qty{1.4e15}{rad/s}$; $f = \qty{2.2e14}{Hz}$; $\lambda = c/f = \qty{1.4}{\micro m}$;
$1/\lambda = \qty{7300}{cm^{-1}}$.

\textbf{11.} Een factor $2.4$ te hoog: de wanden van het model zijn te steil
(de kern met $r^{-12}$ is een grove vervanger voor een covalente binding), maar de
\omterm{def:b1:units-dimensions:oom}{orde van grootte} --- een infraroodtrilling bij $10^{14}$ Hz --- klopt,
en de methode (de kromming van de put) is precies dezelfde als bij
betere potentialen.

\textbf{12.} $\tfrac12 kA^2 \approx k_BT$: $A = \sqrt{2k_BT/k} = \sqrt{8.3\times
10^{-21}/3.1\times10^3} = \qty{1.6e-12}{m} = 0.013\sigma$: de binding
trilt nauwelijks.

\textbf{13.} $v_{\max} = A\omega_0 = 1.6\times10^{-12} \times 1.4\times10^{15}
= \qty{2.2}{km/s}$.

\textbf{14.} $-\epsilon < E_m < 0$: gesloten krommen rond $(\sigma, 0)$ ---
trilling, gebonden; $E_m > 0$: open krommen --- de atomen gaan uiteen;
$E_m = 0$: de separatrix, een atoom dat het oneindige net kan bereiken.

\textbf{15.} $u^2 - 2u = -\tfrac12$ met $u = (\sigma/r)^6$: $u = 1 \pm 1/\sqrt2
= 1.71$ of $0.293$; $r = \sigma u^{-1/6}$: $0.915\sigma$ en $1.23\sigma$.

\textbf{16.} Midden $1.07\sigma > \sigma$: de buitenwand is zachter dan
de binnenwand, zodat het atoom meer ruimte buiten doorbrengt; de over de tijd
gemiddelde afstand groeit met de energie --- verwarmde bindingen worden langer, vaste stoffen zetten uit.

\textbf{17.} Langer: de terugdrijvende kracht aan de zachte buitenkant is
zwakker dan de harmonische benadering aanneemt, zodat de terugkeer
meer tijd kost.

\textbf{18.} De halve putdiepte is $\epsilon/2 = \qty{2.2}{eV} = \qty{3.5e-19}{J}$;
als \omterm{thm:b1:work-and-energy:ke}{kinetische energie} in $\sigma$ geeft dat
$v = \sqrt{2 \times 3.5\times10^{-19}/1.67\times10^{-27}}
= \qty{2.0e4}{m/s}$.

\textbf{19.} Behoud van energie: $E_m > 0$ blijft $> 0$, de kromme is
open, de atomen vliegen na één ontmoeting uiteen. Een derde lichaam (een ander
molecuul, een wand, een foton) moet het overschot afvoeren.

\textbf{20.} $W = E_p(2\sigma) - E_p(\sigma) = \epsilon(2^{-12} - 2 \times 2^{-6})
+ \epsilon = 0.97\epsilon = \qty{4.3}{eV}$, positief: de aantrekking trekt
het atoom naar binnen, dus drijvend.

\textbf{21.} $W_{\text{ext}} = E_p(1.5\sigma) - E_p(\sigma) = \epsilon(0.0077 -
0.176) + \epsilon = 0.83\epsilon = \qty{3.7}{eV}$.

\textbf{22.} $(1/0.9)^{13} = 3.93$, $(1/0.9)^7 = 2.09$: $F = +22\epsilon/\sigma
= \qty{+120}{nN}$, tien keer de aantrekking bij $1.2\sigma$: een binding
$10\%$ samendrukken kost veel meer dan haar $20\%$ uitrekken.

\textbf{23.} $T = \epsilon/k_B = 7.0\times10^{-19}/1.38\times10^{-23} =
\qty{5.1e4}{K}$. Werkelijke gassen dissociëren ver daaronder, omdat de energieën
verdeeld zijn: een deel van de moleculen draagt vele malen $k_BT$
(\cref{ch:b1:kinetic-theory}).

\textbf{24.} Blijven gelden: het evenwicht in het minimum, $k = E_p''$ in het
minimum, dissociatie-energie $=$ diepte, \omterm{def:b1:work-and-energy:turning}{keerpunten} uit $E_p = E_m$,
asymmetrie $\Rightarrow$ uitzetting. Veranderen: de waarden van $k$, van de
frequentie en van de \omterm{def:b1:work-and-energy:turning}{keerpunten}.

\textbf{25.} Evenwicht: $E_p' = 0$; trilling: $\omega_0^2 = E_p''/m$;
ontsnappen: $E_m$ vergeleken met de rand van de put.
\end{solution}
